\setlength{\footskip}{8mm}

\renewcommand{\thefigure}{A.\arabic{figure}}
\renewcommand{\thetable}{A.\arabic{table}}

\begin{center}
{\large \bf Appendix A\\ \vskip 1em Kinematic Tree and BFS Scan Order\\ \vskip 1em}
\vskip 1em
\end{center}

\FloatBarrier

\singlespace

This appendix lists the 17-joint Human3.6M skeleton used throughout the thesis
(Section~\ref{sec:method-problem}) and the breadth-first search (BFS) scan order
derived from it (Section~\ref{sec:method-bfs}).

\vskip 1em
{\bf Joints.} The 17 joints are indexed as: 0 pelvis, 1 right hip, 2 right knee,
3 right ankle, 4 left hip, 5 left knee, 6 left ankle, 7 spine, 8 thorax, 9 neck,
10 head, 11 left shoulder, 12 left elbow, 13 left wrist, 14 right shoulder,
15 right elbow, 16 right wrist.

\vskip 1em
{\bf Bones (kinematic-tree edges).} The 16 parent-child bones, rooted at the pelvis
(joint 0), grouped by chain:

\begin{itemize}
  \item Right leg: $0\rightarrow1\rightarrow2\rightarrow3$
  \item Left leg: $0\rightarrow4\rightarrow5\rightarrow6$
  \item Spine and head: $0\rightarrow7\rightarrow8\rightarrow9\rightarrow10$
  \item Left arm: $8\rightarrow11\rightarrow12\rightarrow13$
  \item Right arm: $8\rightarrow14\rightarrow15\rightarrow16$
\end{itemize}

\vskip 1em
{\bf BFS scan order.} Visiting joints outward from the pelvis, level by level, and
breaking ties by ascending joint index, gives the order
\[
\pi = [\,0,\ 1,\ 4,\ 7,\ 2,\ 5,\ 8,\ 3,\ 6,\ 9,\ 11,\ 14,\ 10,\ 12,\ 15,\ 13,\ 16\,].
\]
The BiSTSSM layer reindexes the joint axis by $\pi$ before the spatial scan and applies
the inverse permutation after the cross-merge, as described in
Section~\ref{sec:method-bfs}.

\FloatBarrier
